\documentclass[12pt]{ctexart}
\usepackage{geometry}
\usepackage{graphicx}
\usepackage{amsmath}
\usepackage{amssymb}

\usepackage{cancel}
\usepackage{hyperref}
\usepackage[svgnames, x11names]{xcolor}

\newcommand{\tG}{\textcolor{DarkGreen}}
\definecolor{graygreen}{RGB}{102,152,126}
\newcommand{\sG}[1]{\textcolor{graygreen}{\sout{#1}}}
\newcommand{\cG}[1]{\textcolor{DarkGreen}{\cancel{#1}}}

\geometry{margin=1in}

\title{一个示例研究论文}
\author{张三}
\date{\today}

\begin{document}

\maketitle

\begin{abstract}
这是一个研究论文的示例摘要。它演示了\texttt{changes}包用于跟踪修订的用法。新的发现表明这种方法对协作写作有效。
\end{abstract}

\section{引言}
\label{sec:intro}

The quick brown fox jumps over the lazy dog. 这是一个段落，演示文本处理。

\subsection{背景}
\label{subsec:background}

之前的工作已经在该领域完成~\cite{ref2021}。重要的贡献包括：
\begin{itemize}
    \item 第一项贡献
    \item 第二项贡献
    \item 第三项贡献
\end{itemize}

\section{方法}
\label{sec:method}

我们使用了以下方程：
\begin{equation}
    E = mc^2
    \label{eq:energy}
\end{equation}

旧方法产生更好的结果。

\subsection{数据收集}
数据来自多个来源。条目已被删除。

\section{结果}
\label{sec:results}

我们的结果如图~\ref{fig:example}所示。数据表明：
\begin{enumerate}
    \item 第一项结果
    \item 第二项结果
    \item 第三项结果
\end{enumerate}

\begin{figure}[h]
    \centering
    % \includegraphics[width=0.5\textwidth]{example.png}
    \caption{示例图}
    \label{fig:example}
\end{figure}

\section{讨论}
\label{sec:discussion}

我们在此讨论我们的发现。不再有效。新的解释提供了更好的见解。

\subsection{该小节只有一段内容}

理论博大精深
每个理论都很厉害 很厉害 厉害。
我学习这些理论
但是好难

\section{结论}
\label{sec:conclusion}

总之，本文演示了\texttt{changes}包的有效性。未来的工作将探索额外的应用。

\bibliographystyle{plain}
\bibliography{references}

\end{document}